\unit{Exposé 236 -- Picard schemes: general properties}
\sid{236}
\src{221}
\seminar{14th year, 1961/62, no. 236\\May 1962}
\paperhead{DESCENT TECHNIQUE AND EXISTENCE THEOREMS IN ALGEBRAIC GEOMETRY\\
VI. PICARD SCHEMES: GENERAL PROPERTIES}%
  {by Alexander GROTHENDIECK}

\fsection{0}{Complements to the preceding exposé ([3], V)\textsuperscript{(1)}}

There has been some progress on the existence questions for Picard preschemes
raised in [3], V:

\medskip
\noindent a. (Mumford). It is not true in general that when
$f:X\longrightarrow S$ is a projective separable morphism (= flat with
separable fibers), the Picard prescheme $\operatorname{Pic}_{X/S}$ exists,
even if the fibers of $f$ have dimension 1 and $S$ is the spectrum of a
complete discrete valuation ring. An example is obtained by taking
$S=\operatorname{Spec}\mathbf R[[t]]$ and $X$ to be the subscheme of
$\mathbf P^2_S$ defined by the equation (where $x,y,z$ are homogeneous
variables)
\[
x^2+y^2=tz^2.
\]
This is a conic degenerating geometrically into two intersecting lines, while
the special fiber over $\mathbf R$ is nevertheless irreducible (it is given by
$x^2+y^2=0$ over $\mathbf R$). One sees readily that after the étale extension
$S'\longrightarrow S$, with $S'=\operatorname{Spec}\mathbf C[[t]]$, the Picard
prescheme of $X'/S'$ exists and has an explicit description as a disjoint union
of copies of $\widetilde S'$, where $\widetilde S'$ is obtained from $S'$ by
doubling the origin infinitely many times. The descent datum on
$\operatorname{Pic}_{X'/S'}$ for $S'\longrightarrow S$ (given here by the
action of the Galois group $\mathbf Z/2\mathbf Z$ of $S'$ over $S$) is readily
seen not to be effective: the group interchanges certain doubled points, so
that some orbits are not contained in an affine open subset. Mumford can,
however, show that if $f:X\longrightarrow S$ is a projective separable
morphism such that, for every $s\in S$, the irreducible components of $X_s$
are geometrically irreducible over $k(s)$, then
$\operatorname{Pic}_{X/S}$ exists. His proof relies on a refinement of his
quotient theorem; cf. [7]. Notice, on the other hand, that without any
hypothesis on the irreducible components of the fibers $X_s$, it remains
possible that the scheme $\operatorname{Pic}^{\tau}_{X/S}$ (to be introduced
below) nevertheless exists.

\src{222}
\noindent b. (Murre). Let $X$ be a proper scheme over a field. Then
$\operatorname{Pic}_{X/k}$ exists. The proof partly follows Chevalley's proof
[1] and makes essential use of the group structure of the Picard functor.

For some complements concerning the theory of Picard schemes, especially in
relation to abelian schemes, it will be useful to consult [7]. Finally, a
notable gap in the present exposé is the absence of “equivalence criteria”
permitting comparison of the Picard scheme of a projective scheme with those
of its hyperplane sections. The key theorems for developing such criteria are
in [5], to which the existence theorems for Picard schemes must be added.

\fsection{1}{Topological properties of commutative group preschemes}

Let first $k$ be a field and $G$ a group prescheme over $k$. Since the identity
element $e$, being rational over $k$, is necessarily closed, it follows at once
that the diagonal of $G\times_kG$ is closed, and hence $G$ is separated: every
group prescheme over a field is separated. Denote by $G^\circ$ the connected
component of the identity element $e$. Since $e$ is rational over $k$,
$G^\circ$ is in fact geometrically connected; that is, formation of $G^\circ$
is compatible with extension of the base field. It also follows that
$G^\circ$ is stable under multiplication as a set, and when $G$ is locally
noetherian, so that $G^\circ$ is open, it may be regarded as an open subgroup
of $G$. Henceforth suppose that $G$ is locally of finite type over $k$. Then
$G^\circ$ is geometrically irreducible and of finite type over $k$. Indeed,
one may suppose $k$ algebraically closed, and then also $G$ reduced (for
$G_{\mathrm{red}}$ is then a subgroup of $G$, since
$G_{\mathrm{red}}\times_kG_{\mathrm{red}}$ is again reduced), hence smooth
over $k$ on a nonempty open subset and therefore everywhere, by translating
that open subset. Thus $G$ is locally irreducible, so its irreducible
components coincide with its connected components, and $G^\circ$ is
irreducible. Let $U$ be an affine neighborhood of $e$ in $G^\circ$. Using the
irreducibility of $G^\circ$, one sees at once that $U\cdot U=G^\circ$, proving
that $G^\circ$ is quasi-compact and hence of finite type over $k$.

For simplicity suppose that $G$ is commutative. For every integer $n>0$, let
$G^{(n)}$ be the inverse image of $G^\circ$ under the $n$th-power
homomorphism $\varphi_n$ in $G$; thus $G^{(n)}$ is an open subgroup of $G$.
Put
\[
G^\tau=\bigcup_nG^{(n)}.
\]

\src{223}
\[
G^\sigma=\bigcup_{(n,p)=1}G^{(n)},
\qquad
G^\rho=\bigcup_hG^{(p^h)},
\]
where $p$ is the characteristic exponent of the field $k$. These are open
subgroups of $G$, satisfying
\[
G^\sigma\cap G^\rho=G^\circ,
\qquad
G^\sigma G^\rho=G^\tau.
\]

\result{REMARK}{} One may construct the quotient group scheme
$G/G^\circ=N$ (cf. [3], IV), and then define $G^\tau$, $G^\sigma$, and
$G^\rho$ as the inverse images in $G$ of the torsion subgroup of $N$
(respectively the natural complement of its $p$-primary component, the sum of
the $q$-primary components for primes $q\neq p$, respectively its $p$-primary
component).
Notice in this connection that $N$ is a discrete separable group scheme over
$k$. Thus, after choosing an algebraic closure $\overline k$ of $k$, with
Galois group $\pi$, it is identified with an ordinary discrete group on which
$\pi$ acts by automorphisms. This gives the evident interpretation of the
construction of the torsion subgroup and of its decomposition into
$q$-primary components. When $G$ is the Picard scheme of a proper scheme $X$
over $k$, $N$ may be called the (reduced) \emph{Néron--Severi scheme} of $X$
over $k$. When $G^\circ_{\mathrm{red}}$ is a group subscheme of $G^\circ$,
which occurs in particular whenever $k$ is perfect or $G^\circ$ is proper over
$k$ (for example, when $X$ is geometrically normal), one should also introduce
the quotient $N'=G/G^\circ_{\mathrm{red}}$, which tends to behave better than
$N$ under specialization, that is, when $X$ varies in a family of algebraic
schemes.

Now let $S$ be a locally noetherian prescheme, and let $G$ be a group
prescheme over $S$, locally of finite type over $S$. We do not assume $G$ to
be of finite type over $S$, or separated over $S$. Put
\[
G^\circ=\bigcup_{s\in S}(G_s)^\circ,
\]
and, when $G$ is commutative,
\[
G^\tau=\bigcup_{s\in S}(G_s)^\tau,
\qquad
G^\sigma=\bigcup_{s\in S}(G_s)^\sigma,
\qquad
G^\rho=\bigcup_{s\in S}(G_s)^\rho.
\]

\src{224}
These are subsets of $G$ stable under multiplication in $G$, which plainly
does not imply that they can be defined by group subpreschemes of $G$. In
particular, it seems that in general there is no group subprescheme of $G$
whose underlying set is $G^\circ$. Of course, if one of these sets is open,
then with its induced structure it is an open group subprescheme of $G$. We
shall see that this is always so for $G^\tau$. Thus, from the standpoint of
representable functors, and especially from the standpoint of specialization,
numerical equivalence behaves more satisfactorily than algebraic equivalence.
Here are the main general properties of the sets just defined:

\result{THEOREM}{1.1} The sets $G^\circ$, $G^\tau$, $G^\sigma$, and $G^\rho$
are locally constructible. Moreover:

\noindent (i) $G^\circ$ is quasi-compact over $S$. When the $G_s^\circ$ are
proper and $G$ is separated over $S$, $G^\circ$ is proper over $S$ and hence
closed in $G$.

\noindent (ii) $G^\tau$ is open. If $G^\circ$ is closed, then so is $G^\tau$.

\noindent (iii) If $G^\circ$ is closed, then so is $G^\sigma$, provided one is
in equal characteristic, that is, all residue fields of $S$ have the same
characteristic. If $G^\circ$ is closed and $G\longrightarrow S$ is universally
open at the points of $G^\sigma$ (cf. Corollary 1.5 below), then
$G^\sigma\longrightarrow S$ is universally open.

\noindent (iv) If $G^\circ$ is closed, then so is $G^\rho$. Suppose one is in
equal characteristic and, for every integer $n>0$ with $(n,p)=1$, the
$n$th-power homomorphism in $G$ is open. Then $G^\rho$ is open.

We give indications for the proof. The local constructibility of $G^\circ$ and
its quasi-compactness over $S$ are contained in the following lemma:

\result{LEMMA}{1.2} Suppose $S\neq\varnothing$. Then there are a nonempty open
subset $U$ of $S$, a finite-type group scheme $H$ over $U$ with connected
fibers, and a homomorphism of group schemes $H\longrightarrow G|U$ that is an
open immersion with image $G^\circ|U$.

To prove the lemma, one may suppose $S$ irreducible and let $\eta$ be its
generic point. After the base change
$S'=\operatorname{Spec}(\mathcal O_{S,\eta})\longrightarrow S$, one obtains a
group scheme $G'$ over the local artinian ring $\mathcal O_{S,\eta}=A$ with an
open finite-type group subscheme $G'{}^\circ$ over $A$, as stated above. It
therefore comes from a finite-type group scheme $H$ over an open neighborhood
$U$ of $\eta$, and the canonical immersion
$G'{}^\circ\longrightarrow G'$ comes from an

\src{225}
open immersion $H\longrightarrow G|U$, which is a group-scheme homomorphism
after $U$ is made sufficiently small. If $U$ is sufficiently small, the fibers
of $H$ are connected and all have the same dimension, equal to that of the
fibers of $G$. It follows that for every $s\in U$ the image of $H_s$ in $G_s$
is exactly $G_s^\circ$ (after again shrinking $U$), proving 1.2.

The second assertion of (i) is contained in the following lemma, applied to a
quasi-compact open neighborhood of $G^\circ$ in $G$:

\result{LEMMA}{1.3} Let $X$ be a finite-type separated prescheme over a locally
noetherian $S$, let $g$ be a section of $X$ over $S$, and let $X^\circ$ be the
union, in the $X_s$, of the connected components containing $g(s)$. If
$s\in S$ and $X_s^\circ$ is proper over $k(s)$, then there is an open
neighborhood $U$ of $s$ such that $X^\circ|U$ is proper over $U$, and hence
closed in $X|U$.

By faithfully flat descent on the base, one reduces to the case in which $S$
is the spectrum of a complete local ring and $s$ is its closed point. Applying
[2], III, 5.5.1, one sees that $X$ is the disjoint union of two open subsets
$X'$ and $X''$, the first proper over $S$ and satisfying
$X'_s=X_s^\circ$. This reduces to the case $X=X'$, that is, $X$ proper over
$S$. One then uses a standard argument based on the valuative criterion for
properness of a subset (omitted from [2], Chapter II).

We now prove that $G^\tau$ is open. Equivalently, since formation of $G^\tau$
(as well as of $G^\circ$, $G^\sigma$, and $G^\rho$) commutes with extension of
the base, for every section $g$ of $G$ over $S$, the set $g^{-1}(G^\tau)$ is
open. This means two things:

\medskip
\noindent a. If $y\in S$ and $g(y)\in G_y^\tau$, then
$g(y')\in G_{y'}^\tau$ for every $y'\in\overline{\{y\}}$ sufficiently near
$y$.

\medskip
\noindent b. If $y'\in S$ and $g(y')\in G_{y'}^\tau$, then
$g(y)\in G_y^\tau$ for every generalization $y$ of $y'$.

For (a), some integer $n>0$ satisfies $g^n(y)\in G_y^\circ$. Since $G^\circ$
is constructible, $(g^n)^{-1}(G^\circ)$ is constructible, and hence
$g^n(y')\in G_{y'}^\circ$ for $y'\in\overline{\{y\}}$ sufficiently near $y$.
For (b), the powers $g(y')^n=g^n(y')$ remain in a quasi-compact open subset of
$G_{y'}$, since they lie in finitely many classes modulo $G_{y'}^\circ$.
Consequently there is a quasi-compact open subset $U$ of $G$ containing the
$g^n(y')$ and therefore also their generalizations $g^n(y)$. Thus the powers
of $g(y)$ remain in a quasi-compact open subset of $G_y$, which readily implies
$g(y)\in G_y^\tau$.

\src{226}
Suppose $G^\circ$ is closed; we prove the same for $G^\tau$. Since $G^\tau$ is
already known to be open and hence locally constructible, it remains to prove
that it is stable under specialization. This follows because it is a union of
closed sets, namely the inverse images of the closed set $G^\circ$ under the
$n$th-power homomorphisms $\varphi_n$.

The same argument proves that $G^\sigma$ and $G^\rho$ are closed if $G^\circ$
is closed (in the first case subject to equal characteristic), once it is shown
that $G^\sigma$ and $G^\rho$ are locally constructible. For $x\in G^\tau$, let
$\nu(x)$ be the smallest integer $n>0$ such that the $n$th-power homomorphism
$\varphi_n$ sends $x$ into $G^\circ$. Then $G^\sigma$ (respectively $G^\rho$)
consists of those $x\in G$\footnote{The printed text reads “$X$,” which is not
defined in this passage; all four loci under consideration are subsets of
$G$.} for which $\nu(x)$ is prime to $p$ (respectively a power of $p$). The
constructibility assertion now follows from the more precise statement:

\result{LEMMA}{1.4} The function $\nu$ on $G^\tau$ is locally constructible.

Indeed, for every integer $n>0$, the set of $x\in G^\tau$ such that
$\nu(x)=n$ is locally constructible: it is the difference between
$\varphi_n^{-1}(G^\circ)$ and the union of the
$\varphi_d^{-1}(G^\circ)$, where $d$ ranges over the proper divisors of $n$.
Since $G^\circ$ is locally constructible, so are the
$\varphi_d^{-1}(G^\circ)$, and therefore so is the preceding difference.

Suppose now that $G^\circ$ is closed and $G\longrightarrow S$ is universally
open at the points of $G^\sigma$. We prove that
$G^\sigma\longrightarrow S$ is open, that is, it maps a neighborhood in
$G^\sigma$ of $x\in G^\sigma$ onto a neighborhood of $y=f(x)$. Since
$G^\sigma$ is locally constructible, it suffices to prove that for every
generalization $y'$ of $y$ there is a generalization $x'$ of $x$ in
$G^\sigma$ lying over $y'$. Base change reduces this to the case in which $S$
is the spectrum of a discrete valuation ring and $y,y'$ are respectively its
closed and generic points. Since $G\longrightarrow S$ is open at $x$, there is
a generalization $x_1$ of $x$ in $G$ above $y'$. After a further base change,
one may also suppose that there is a section $g$ of $G$ over $S$ with
$x=g(y)$. If $k(y')$ has characteristic zero, take any generalization $x'$ of
$x$ in $G$ above $y'$. It belongs to $G^\tau$ because $G^\tau$ is open, hence
to $G^\sigma$ because $G_{y'}^\sigma=G_{y'}^\tau$. If $k(y')$ has
characteristic $p>0$, write
\[
\nu(g(y'))=p^hm\qquad\text{with }(m,p)=1;
\]

\src{227}
choose integers $a$ and $b$ such that $ap^h+bm=1$, and set
\[
g_1=g^{ap^h},\qquad g_2=g^{bm},\qquad\text{so that}\qquad g=g_1g_2.
\]
By construction, $g_1(y')\in G^\sigma$ and $g_2(y')\in G^\rho$. Since
$G^\circ$ is closed, it follows that $g_1(y)\in G^\sigma$ and
$g_2(y)\in G^\rho$. But
\[
g(y)=g_1(y)g_2(y)\in G^\sigma,
\]
so $g_2(y)\in G^\sigma$, and therefore
\[
g_2(y)\in G_y^\sigma\cap G_y^\rho=G_y^\circ.
\]
The hypothesis, together with the fact that $S$ is the spectrum of a discrete
valuation ring, implies that $G-(G_y-G_y^\circ)$ is an open subset of $G$ on
which $G\longrightarrow S$ induces an open morphism. Thus a “quasi-section”
passes through every point of $G_y^\circ$. After making a further extension of
the base $S$, one may consequently suppose that there is a section $g'_2$ of
$G^\circ$ over $S$ such that $g'_2(y)=g_2(y)$. Set $g'=g_1g'_2$. Then, by
construction, $g'(y)=g(y)=x$, and
\[
g'(y')=g_1(y')g'_2(y')\in G^\sigma.
\]
Hence $g'(y')=x'$\footnote{The printed text reads “$g(y')=x'$.” The section
constructed in the preceding sentence, whose generic fiber is being computed,
is $g'$.} is a generalization of $x$ in $G^\sigma$ lying above $y'$. This
proves (iii).

It remains to prove the last assertion, (iv). We must show that if
$x\in G^\rho$, then every generalization $x'$ of $x$ belongs to $G^\rho$. By
taking images under $\varphi_{p^h}$ for a suitable $h$\footnote{The printed
text reads $\varphi_{n^h}$. For $x\in G^\rho$, the definition on p.~223 gives
a power of the characteristic exponent $p$ that sends $x$ into $G^\circ$.},
we may suppose that in fact $x\in G^\circ$. For every integer $n$ prime to the
characteristic, $x$ then lies in the image of $\varphi_n$ (because the
$n$th-power map in a connected group of finite type over a field of
characteristic prime to $n$ is surjective). Since $\varphi_n$ is open, $x'$
also lies in the image of $\varphi_n$. More precisely, let $U$ be a
quasi-compact open subset of $G^\tau$ containing $G_y^\circ$. Then
$x'\in\varphi_n(U)$ for every $n$ prime to $p$. Taking $n$ to be a common
multiple of the factors prime to $p$ of the $\nu(z)$, for $z\in U$, gives
$x'\in G^\rho$.

Assertion (1.1), (iii), is completed by the following:

\result{COROLLARY}{1.5} Let $n>1$ be an integer such that the $n$th-power
homomorphism $\varphi_n:G\longrightarrow G$ is universally open (for example,
étale), put
\[
G^{(n)}=\varphi_n^{-1}(G^\circ),
\]
and suppose that, in every geometric fiber $G_{\bar s}$, the subset

\src{228}
$\bigcup_{h\geq0}|{}_{n^h}G_{\bar s}|$ is dense in
$\bigcup_{h\geq0}G_{\bar s}^{(n^h)}$. Then
$G\longrightarrow S$ is universally open at the points of the
$G^{(n^h)}$. In particular, if the $G_s^\circ$ contain no additive component
and if, for every integer $n>1$, the homomorphism $\varphi_n$ is universally
open at the points of residue characteristic prime to $n$\footnote{The printed
text reads “prime to $p$.” In this statement quantified by $n$, the étaleness
criterion and the argument using the kernels ${}_nG$ require residue
characteristic prime to $n$.}, then $G\longrightarrow S$ is universally open
at the points of $G^\sigma$. Hence, by (1.1), (iii), if $G^\circ$ is closed,
$G^\sigma\longrightarrow S$ is universally open. In these cases,
$s\longmapsto\dim G_s$ is locally constant on $S$.

Indeed, the hypothesis implies that the kernel ${}_nG$ of $\varphi_n$ is
universally open over $S$. Hence $G\longrightarrow S$ is universally open at
the points of ${}_nG$, and therefore, on replacing $n$ by $n^h$, at the points
of ${}_{n^h}G$. The density hypothesis says exactly that, fiber by fiber, the points of order
a power of $n$ are dense in the union of the $G^{(n^h)}$. For the particular
assertion, absence of additive components implies this density for integers
$n$ invertible in the residue field. It follows readily that $G\longrightarrow S$ is universally open
at every point of this union, in particular along the identity section, and
hence that $s\longmapsto\dim G_s$ is locally constant.

\result{REMARK}{1.6} Recall that a morphism $X\longrightarrow S$ is called
universally open if it maps every open subset to an open subset and retains
this admirable property after every base change $S'\longrightarrow S$. This
also means (when $S$ is locally noetherian and $X\longrightarrow S$ is locally
of finite type) that every irreducible component of $X$ dominates $S$, and
that this property is preserved after every base change
$S'\longrightarrow S$. In both assertions it is enough (under the finiteness
hypotheses above) to test base changes $S'\longrightarrow S$ in which $S'$ is
the spectrum of a discrete valuation ring (complete with algebraically closed
residue field, if desired~...). The definition extends in the evident way to a
subset $Z$ of $X$ (such as the subset $G^\sigma$ of $G$). It is a perfectly
normal phenomenon, even for a projective smooth morphism
$X\longrightarrow S$ with geometrically connected fibers (for example, the
fiber square of the modular family of elliptic curves over
$S=\operatorname{Spec}\mathbf C[j]$), that $\operatorname{Pic}_{X/S}$ is not
universally open over $S$: an irreducible component of
$\operatorname{Pic}_{X/S}$ may lie entirely above a single point of $S$. This
is connected with the fact that the rank of the Néron--Severi group of the
fibers of $X/S$ can jump upward (“complex multiplication” phenomena). On the
other hand, (1.5) assures us that in favorable cases
$\operatorname{Pic}_{X/S}^\sigma$ (and most often even
$\operatorname{Pic}_{X/S}^\tau$, it seems) is universally open over $S$.

\src{229}
Finally, here is a useful case in which, exceptionally, $G^\circ$ consents to
be open:

\result{COROLLARY}{1.7} Suppose that $G\longrightarrow S$ is universally open
at the points of $G^\circ$ (cf. (1.5)), and that the fibers $G_s$ are separable
and hence smooth over $k(s)$ (the latter condition is automatically satisfied
in residue characteristic zero, by a theorem of Cartier). Then $G^\circ$ is
open in $G$. If, in addition, $S$ is reduced, then $G^\circ$ is smooth (in
particular, flat) over $S$.

The first assertion can be sharpened as follows. If, for a given $s\in S$,
$G\longrightarrow S$ is universally open at the points of $G_s^\circ$ and
$G_s^\circ$ is separable over $k(s)$, then $G^\circ$ is a neighborhood of
$G_s^\circ$ (and the hypothesis made at $s$ will moreover remain valid at
neighboring points). This statement, which is independent of any group
structure, will be proved in [2], IV, §7. The last assertion in (1.7), also
independent of any group structure or any question about a connected
component, is a special case of a flatness criterion given in [2], IV, §5, which
more generally implies:

\result{COROLLARY}{1.8} Let $U$ be an open subset of a fiber $G_s$ such that
$G\longrightarrow S$ is universally open at the points of $U$ (cf. (1.5)). If
$G_s$ is separable over $k(s)$ and $S$ is reduced at $s$, then $G$ is flat over
$S$ at the points of $U$ (and thus, in this case, smooth over $S$ at the points
of $U$).

\result{REMARKS}{1.9} I do not know whether, when $G$ is separated over $S$,
$G^\circ$ or $G^\tau$ is always closed in $G$; this seems unlikely. There are,
at any rate, obvious counterexamples if the separatedness hypothesis is
dropped. For example, take the affine line with its origin doubled countably
infinitely many times. This gives a group scheme $G$ over the affine line all
of whose fibers are reduced to the identity group except for one, whose fiber
is $\mathbf Z$. This group scheme is an open subgroup scheme, namely the
closure of the identity section, in the Picard prescheme of the $S$-scheme $X$
corresponding to a family of conics degenerating into two intersecting lines.

Moreover, even starting with a finite flat group scheme $G$ over the spectrum
$S$ of a discrete valuation ring $V$, so that $G^\circ$ is reduced to the
identity section and hence closed, various conclusions become false if some
hypotheses are dropped. Suppose that $V$ has equal characteristic $p>0$, and
let $G$ be the kernel of the homomorphism
$\mathbf G_a\longrightarrow\mathbf G_a$ defined by the homomorphism of
functors
\[
f\longmapsto f^p-tf,
\]
where $t$ is a uniformizer of $V$.

\src{230}
(Recall that, by definition, the “additive group” $\mathbf G_a$ over $S$
represents the functor $S'\longmapsto\Gamma(S',\mathcal O_{S'})$.) Thus, as an
$S$-scheme,
\[
G=\operatorname{Spec}V[x]/(x^p-tx)
\]
is the union of $p$ intersecting “lines,” representing a group
$\mathbf Z/p\mathbf Z$ that degenerates into an infinitesimal group of additive
type. In this example, $G^\circ=G^\sigma$ (and it is one of the $p$ lines), so
it is not open in $G$. In mixed characteristic with residue characteristic
$p>0$, start with the group scheme $H=\boldsymbol\mu_p$, the kernel of the
$p$th-power map in $\mathbf G_m$. Thus
\[
H=\operatorname{Spec}V[x]/(x^p-1);
\]
this is again the union of $p$ intersecting lines, representing a group
$\mathbf Z/p\mathbf Z$ (in characteristic zero) that degenerates into an
infinitesimal group of multiplicative type in characteristic $p$. Let $H'$ be
the “constant group scheme” defined by the ordinary finite group
$\mathbf Z/p\mathbf Z$, that is, the disjoint union of $\mathbf Z/p\mathbf Z$
copies of $S$, or equivalently $\operatorname{Spec}V^{\mathbf Z/p\mathbf Z}$.
Then $G=H\times_S H'$ describes a group $(\mathbf Z/p\mathbf Z)^2$ in
characteristic zero, degenerating in characteristic $p$ into an infinitesimal
group times $\mathbf Z/p\mathbf Z$. Here $G^\rho$ is the union of the identity
section and the special fiber, so $G^\rho$ is not open, whereas $G^\sigma$ is
the union of the identity section and the generic fiber, so it is not closed,
contrary to what is asserted in the equal-characteristic cases in (1.1), (iii)
and (iv). These are, of course, phenomena connected with characteristic $p>0$.
The preceding results give:

\result{COROLLARY}{1.10} Suppose that all residue characteristics of $S$ are
zero, so that $G^\tau=G^\sigma$ and $G^\circ=G^\rho$. Then
$G^\tau=G^\sigma$ is open; it is even open and closed if $G$ is separated over
$S$ and the $G_s^\circ$ are proper. Under the same hypothesis,
$G^\circ=G^\rho$ is proper over $S$ and hence closed in $G$. Finally, suppose
that for every integer $n>1$ the $n$th-power homomorphism in $G$ is universally
open. Then $G^\circ$ is open; and if, moreover, the $G_s^\circ$ have no
additive component (for example, if the $G_s^\circ$ are proper, as above),
then $G^\tau\longrightarrow S$ is universally open, and is even smooth if $S$
is reduced.

We finally record the following easy result:

\result{PROPOSITION}{1.11} There is an open subset $U$ of $S$ such that the
set $G'$ of points of $G$ at which $G$ is smooth (respectively flat) over $S$
is the underlying set of an open subgroup scheme induced from $G|U$.
Furthermore, every section of $G$ over $U$ is a section of $G'$ over $U$.

\src{231}
\result{COROLLARY}{1.12} If $G$ is smooth (respectively flat) over $S$ at the
points of the identity section, then it is so at the points of every section
of $G$ over $S$ and at every point of $G^\circ$. If, moreover, for every
integer $n>0$ the $n$th-power homomorphism
$\varphi_n:G\longrightarrow G$ is étale at the points of characteristic prime
to $n$, then $G$ is smooth (respectively flat) over $S$ at every point of
$G^\sigma$.

\fsection{2}{Application to local properties of Picard schemes}

\result{THEOREM}{2.1}

\noindent(i) Let $f:X\longrightarrow S$ be a proper smooth morphism, and
suppose that $\operatorname{Pic}_{X/S}$ exists (for example, suppose $f$ is
projective). Then $\operatorname{Pic}_{X/S}$ is separated over $S$, and every
closed subset $Z$ of $\operatorname{Pic}_{X/S}$ that is of finite type over
$S$ is proper over $S$.

\noindent(ii) Let $X$ be a proper geometrically normal prescheme over a field
$k$. Then $\operatorname{Pic}_{X/k}^\circ$\footnote{The printed text reads
$\operatorname{Pic}_{X/S}^\circ$, although $S$ does not occur in this
statement over the field $k$.} is proper over $k$.

\begin{proof}
(i) The valuative criteria ([2], II, §7) reduce us to proving the following. Let
$S$ be the spectrum of a complete discrete valuation ring and $U$ the open
subset consisting of the generic point of $S$. Every rational section of
$\operatorname{Pic}_{X/S}$ over $S$, that is, every section over $U$, extends
uniquely to a section over $S$. By the definition of
$\operatorname{Pic}_{X/S}$, this is equivalent to the following statement: for
every invertible Module $\mathcal L$ on $V=f^{-1}(U)$ there is an invertible
Module on $X$ extending $\mathcal L$, unique up to isomorphism. This follows
readily from the description of invertible Modules on $V$, respectively on
$X$, in terms of classes of Cartier divisors. Indeed, the local rings of $X$
are regular ($X$ being smooth over the regular scheme $S$), hence factorial by
Auslander's theorem; thus every divisor on $X$\footnote{The printed text reads
“on $S$.” The factoriality invoked is that of the local rings of $X$.} is a
Cartier divisor. Every divisor on $V$ extends to a divisor on $X$ by taking its
closure.
\end{proof}

\result{REMARK}{2.2} The proof remains valid if one assumes only that $f$ is
flat and that its fibers $X_s$ are locally complete intersections and smooth
over $k(s)$ in codimension $\leq2$. This follows from the theorem proved in
[5]: a noetherian local complete-intersection ring that is regular in
codimension $\leq3$ is factorial (“Samuel's conjecture”). Notice that the

\src{232}
result becomes false if “codimension $\leq2$” is replaced by “codimension
$\leq1$,” that is, by the hypothesis “normal,” as is seen from a family of
nonsingular quadrics degenerating into a quadratic cone.

(ii) Chow's lemma reduces us to the case in which $X$ is projective and hence
embedded in $\mathbf P_k^n$; we may suppose that $X$ is connected. If
$\dim X=1$, then $X$ is smooth over $k$, and (i) applies. If $\dim X\geq2$, a
variant of the known “equivalence criteria” shows that there are finitely many
curves $Y_i$, smooth on $X$ (obtained as intersections of $X$ with suitable
linear subspaces of $\mathbf P_k^n$), such that
\[
\operatorname{Pic}_{X/k}^\tau\longrightarrow
\prod_i\operatorname{Pic}_{Y_i/k}^\tau
\]
is a monomorphism, and hence induces a monomorphism on connected components.
The connected component of the right-hand side is proper over $k$ by the
preceding result, and the homomorphism of group schemes in question is
necessarily a closed immersion. It follows that
$\operatorname{Pic}_{X/k}^\circ$ is also proper over $k$. The delicate
equivalence criteria can be avoided by using the structure of commutative
algebraic groups over an algebraically closed field (due to Chevalley--Borel).
One is reduced to proving that every morphism from the affine line with the
origin removed\footnote{The printed text reads “from the affine line.” The
linked erratum, p.~303, prescribes “affine line with the origin removed.”} into
$\operatorname{Pic}_{X/k}^\tau$ is constant. This is equivalent to saying that
every invertible Module on $X[t,t^{-1}]$\footnote{The printed text reads
$X[t]$; the linked erratum, p.~303, prescribes $X[t,t^{-1}]$.} comes from an
invertible Module on $X$. This well-known elementary result does not even use
the properness of $X$ over $k$ (the normality of $X$ immediately reduces one
to the case in which $X$ is regular).

\result{COROLLARY}{2.3} Let $f:X\longrightarrow S$ be a proper normal morphism
(that is, flat with geometrically normal fibers), and suppose that
$\operatorname{Pic}_{X/S}$ exists. Then $\operatorname{Pic}_{X/S}^\circ$ is
proper over $S$ and hence closed in $\operatorname{Pic}_{X/S}$. Moreover,
$\operatorname{Pic}_{X/S}^\tau$ and $\operatorname{Pic}_{X/S}^\rho$ are also
closed, as is $\operatorname{Pic}_{X/S}^\sigma$ in “equal characteristic.”

Apply (1.1) and (2.1), (ii).

\result{COROLLARY}{2.4} Let $f:X\longrightarrow S$ be a proper smooth morphism
such that $\operatorname{Pic}_{X/S}$ exists and is a disjoint union of schemes
$P^{(i)}$ of finite type over $S$ (cf. [3], V, 4.1). Then each $P^{(i)}$ is
proper over $S$.

This follows from (2.1), (i). As noted in (2.2), the result can be generalized
by making less restrictive hypotheses on the fibers of $f$, but it becomes
false if one merely assumes that $f$ is normal. In that case, I do not know
whether $\operatorname{Pic}_{X/S}^\tau$ is nevertheless proper over $S$, even
assuming it to be of finite type over $S$.

\src{233}
\result{THEOREM}{2.5} Let $f:X\longrightarrow S$ be a proper flat morphism
such that $\operatorname{Pic}_{X/S}$ exists, and for every integer $n$ let
$\varphi_n$ be the $n$th-power homomorphism in this group prescheme. Then
$\varphi_n$ is étale at every point
$x\in\operatorname{Pic}_{X/S}$\footnote{The printed text reads “$x\in X$.”
The map $\varphi_n$ is the power endomorphism of the Picard prescheme.} whose
residue characteristic is prime to $n$.

By the infinitesimal characterization of étale morphisms, this assertion is
equivalent to the following:

\result{LEMMA}{2.6} Suppose that $S$ is the spectrum of a local Artinian ring
$A$ whose maximal ideal $\mathfrak m$ has zero $(\nu+1)$st power, and put
\[
A_{\nu-1}=A/\mathfrak m^\nu,
\qquad X_{\nu-1}=X\otimes_A A_{\nu-1}.
\]
Let $\mathcal L$ be an invertible Module on $X$, and let
$\mathcal L'_{\nu-1}$ be an invertible Module on $X_{\nu-1}$ whose $n$th
tensor power is isomorphic to
\[
\mathcal L_{\nu-1}=\mathcal L\otimes_A A_{\nu-1}.
\]
Then there is an invertible Module $\mathcal L'$ on $X$ whose $n$th tensor
power is isomorphic to $\mathcal L$, provided that $n$ is prime to the residue
characteristic of $k=k(A)$.

We prove the lemma. Put
\[
V=\mathfrak m^\nu=\mathfrak m^\nu/\mathfrak m^{\nu+1};
\]
this is a vector space over $k=k(A)$. First try to extend
$\mathcal L'_{\nu-1}$ to an arbitrary invertible Module $\mathcal L'$ on $X$.
The obstruction lies in
\[
H^2(X_0,\mathcal O_{X_0})\otimes_k V.
\]
The hypothesis on $\mathcal L'_{\nu-1}$ and the extendability of
$\mathcal L_{\nu-1}$ show that $n$ times this obstruction is zero; the
obstruction itself is therefore zero because $n$ is prime to the
characteristic. The ambiguity in the chosen extension lies in
\[
H^1(X_0,\mathcal O_{X_0})\otimes_k V,
\]
and the discrepancy $\xi$ between $\mathcal L'^{\otimes n}$ and $\mathcal L$
lies in the same Module. Correcting $\mathcal L'$ to make this discrepancy zero
amounts to finding an $\eta$ in that Module such that $n\eta=\xi$. This is
again possible because $n$ is prime to the characteristic.

\result{COROLLARY}{} Under the hypotheses of (2.5), suppose in addition that
the Picard schemes of the fibers $X_s$ contain no additive component (for
example, that the $X_s$ are geometrically normal; cf. (2.1), (ii)). Then
$\operatorname{Pic}_{X/S}\longrightarrow S$ is universally open at the points
of $\operatorname{Pic}_{X/S}^\circ$. If
$\operatorname{Pic}_{X/S}^\circ$ is closed (for example, if the $X_s$ are
geometrically normal; cf. (2.3)), then $\operatorname{Pic}_{X/S}^\sigma$ is
itself universally open over $S$. Finally, in equal characteristic,
$\operatorname{Pic}_{X/S}^\rho\longrightarrow S$ is universally open.

Apply (1.5) and (1.1).

\src{234}
\result{COROLLARY}{2.7} Let $f:X\longrightarrow S$\footnote{The printed text
reads “$f:X\longrightarrow Y$,” although both the Picard prescheme and the
following function have base $S$.} be a proper flat morphism such that
$\operatorname{Pic}_{X/S}$ exists. Then the function
\[
s\longmapsto\dim\operatorname{Pic}_{X_s/k(s)}
\]
on $S$ is upper semicontinuous (that is, it can jump upward but not downward),
and is even continuous (that is, locally constant) if the
$\operatorname{Pic}_{X_s/k(s)}$ contain no additive component.

The first assertion is trivial, or nearly so, for every group prescheme
locally of finite type over a locally noetherian base, since it is enough to
look along the identity section. The second assertion follows from (2.5).

\result{REMARK}{2.8} Let $s,s'\in S$, with $s$ a specialization of $s'$. Then
(2.7) amounts to an inequality (respectively an equality) between the
dimensions of the Picard schemes of $X_{s'}$ and of its “specialization” $X_s$.
Even before Picard schemes were constructed, Serre had observed that, for a
smooth morphism $f:X\longrightarrow S$, invariance of the dimension of the
Picards of the $X_s$ is a formal consequence of the specialization theory of
the fundamental group ([4], X) and the classical Kummer relations between the
points of finite order on the Picard variety and the abelianized fundamental
group ([4], XI). In general, let $\alpha,\mu,\lambda$ denote the dimensions of
the abelian, multiplicative, and additive parts of
$\operatorname{Pic}_{X_s/k(s)}^\circ$, and define
$\alpha',\mu',\lambda'$ similarly. The known relations are expressed by the
following inequalities:
\[
\tag{*}\alpha+\mu+\lambda\geq\alpha'+\mu'+\lambda'
\]
(valid provided that $\operatorname{Pic}_{X/S}$ exists, and therefore probably
in all cases). When $\lambda=\lambda'=0$, this inequality becomes an equality,
valid under the same existence hypotheses:
\[
\alpha+\mu=\alpha'+\mu';
\]
on the other hand,
\[
\tag{**}2\alpha+\mu\leq2\alpha'+\mu'
\]

\src{235}
by Serre's argument if the $X_s$ are separable (without even assuming the
existence of $\operatorname{Pic}_{X/S}^\tau$), or if the
${}_n\operatorname{Pic}_{X/S}$ (the kernels of $\varphi_n$ in the Picard
functor) are separated over the open of $S$ where $n$ is invertible (since
they are étale over that open by (2.5)). One
is inclined to conjecture that (*) is always an equality, or at least when the
$X_s$ are separable, and that
\[
\tag{***}\alpha\leq\alpha',\qquad\lambda\geq\lambda'
\]
holds whenever there is a group prescheme, locally of finite type over a
locally noetherian $S$, whose fibers have constant dimension (see (1.3) for a
positive result in this direction).

\result{REMARK}{2.9} In every known case,
$\operatorname{Pic}_{X/S}^\tau$ is universally open over $S$, but one should
probably not let this encourage excessive illusions, even when
$f:X\longrightarrow S$ is smooth. In any event, Mumford has constructed an
example (admittedly with $S$ nonreduced, in fact the spectrum of an Artinian
ring) in which $\operatorname{Pic}_{X/S}^\tau$ is not flat over $S$, by varying
the Igusa surface infinitesimally. The point considered by Mumford lies in
$\operatorname{Pic}_{X/S}^\rho$, and it remains possible, for a smooth
$f:X\longrightarrow S$, that $\operatorname{Pic}_{X/S}^\tau$ is flat over $S$
at the points of $\operatorname{Pic}_{X/S}^\sigma$. The speaker nevertheless
doubts that this is always so, even if one restricts to the points of
$\operatorname{Pic}_{X/S}^\circ$ and assumes that $S$ is the spectrum of a
discrete valuation ring. The question is also connected with the study of
fixed points of an abelian scheme under a finite group of automorphisms, a
situation for which examples seem to be lacking.

It would seem that even when one restricts to a smooth projective morphism
$f:X\longrightarrow S$, the local regularity results for
$\operatorname{Pic}_{X/S}$ stated in this section, together with the
conjectures in (2.8), nearly exhaust what can be said without more specific
hypotheses on the nature of the fibers of $f$. Recall, however, that if the
geometric fibers of $\operatorname{Pic}_{X/S}$ are reduced and have no
additive component, then (1.8) and (2.5) imply that
$\operatorname{Pic}_{X/S}$ is smooth over $S$ at the points of
$\operatorname{Pic}_{X/S}^\sigma$ when $S$ is reduced. The same result holds,
without any hypothesis on $S$, if $f:X\longrightarrow S$ is normal, as we
shall see in (3.5). We record in this connection:

\result{PROPOSITION}{2.10}

\noindent(i) If $\operatorname{Pic}_{X/S}$ is smooth (respectively flat) over
$S$ at the points of the identity section, then it is so at every point of
$\operatorname{Pic}_{X/S}^\sigma$ and at the points of every section of
$\operatorname{Pic}_{X/S}$ over $S$.

\src{236}
\noindent(ii) Let $s\in S$ satisfy
\[
H^2(X_s,\mathcal O_{X_s})=0.
\]
Then there is an open neighborhood $U$ of $s$ such that
$\operatorname{Pic}_{X/S}|U$ is smooth over $U$.

\noindent(iii) If $X$ is a proper scheme over a field, then
\[
\dim\operatorname{Pic}_{X/k}\leq\dim H^1(X,\mathcal O_X),
\]
with equality if and only if $\operatorname{Pic}_{X/k}$ is smooth over $k$;
this is always the case when $k$ has characteristic zero.

Part (i) follows from (2.5) and (1.12), and (ii) follows from the infinitesimal
criterion for smooth morphisms and a familiar obstruction calculation, since
the “semicontinuity theorem” shows that the hypothesis at $s$ remains valid at
neighboring points. Finally, (iii) follows from the fact that
$H^1(X,\mathcal O_X)$ is isomorphic to the Zariski tangent space of
$\operatorname{Pic}_{X/k}$ at the identity. The last assertion is a special
case of a theorem of Cartier saying that a “formal group” in characteristic
zero is formally smooth over $k$.

\fsection{3}{The canonical abelian subscheme of
$\operatorname{Pic}_{X/S}$, and the Albanese scheme}

\result{PROPOSITION}{3.1} Let $k$ be a field and $G$ a commutative group
scheme of finite type over $k$ with no additive component. Then
$G_{\mathrm{red}}^\circ$ is separable over $k$, and consequently is a smooth
subgroup scheme over $k$.

This is trivial when $k$ is perfect, in particular for $G_{\bar k}$, where
$\bar k$ is an algebraic closure of $k$. It therefore suffices to see that
$(G_{\bar k}^\circ)_{\mathrm{red}}$ descends from a subscheme of $G$. The
hypothesis on $G_{\bar k}$ (no additive component)\footnote{The printed text
has a malformed construction equivalent to “The hypothesis ... results.” The
linked erratum, p.~303, inserts the preposition needed for “It follows from the
hypothesis ...”.} readily implies that there
is an integer $m$ such that $(G_{\bar k}^\circ)_{\mathrm{red}}$ is the
scheme-theoretic image of the $m$th-power homomorphism in $G_{\bar k}$. Since
this homomorphism comes from the analogous homomorphism for $G$, the
scheme-theoretic image of the latter gives the required subscheme.

\result{COROLLARY}{3.2} Let $X$ be a geometrically normal proper scheme over $k$ such that
$\operatorname{Pic}_{X/k}$ exists. Then there is an abelian subscheme $A$ of
$\operatorname{Pic}_{X/k}$ whose underlying set is
$\operatorname{Pic}_{X/k}^\circ$.

Indeed, by (2.1), (ii), $\operatorname{Pic}_{X/k}^\circ$ is proper over $k$ and
therefore satisfies the condition of (3.1).

\src{237}
Thus the preceding result shows that in certain cases the classical “Picard
variety” (which in the present theory would be denoted
$(\operatorname{Pic}_{X/k})_{\mathrm{red}}^\circ$) “is defined over $k$,”
without assuming that $k$ is perfect.

Now let $f:X\longrightarrow S$ be a proper flat relative scheme, with
\[
\mathcal O_S\xrightarrow{\sim}f_*(\mathcal O_X)
\]
for simplicity, such that $\operatorname{Pic}_{X/S}$ exists and
$\operatorname{Pic}_{X/S}^\circ$ is proper over $S$. For (3.3), (ii), we also
assume that there is an open subset of $\operatorname{Pic}_{X/S}$ containing
$\operatorname{Pic}_{X/S}^\circ$ that is quasi-projective over $S$. As seen
earlier, this condition holds when $f$ is projective with geometrically
separable irreducible fibers. Recall that an abelian scheme over $S$ is a group
scheme over $S$ that is proper and smooth over $S$ and has geometrically
connected fibers. We ask whether there is a subgroup scheme $A$ of
$\operatorname{Pic}_{X/S}$ which is an abelian scheme and whose underlying set
is $\operatorname{Pic}_{X/S}^\circ$. We have just seen that one always exists
when $S$ is the spectrum of a field. Here is what can be said in the general
case under consideration:

\result{THEOREM}{3.3} Under the preceding hypotheses:

\noindent(i) If there is an abelian subscheme of
$\operatorname{Pic}_{X/S}$ whose underlying set is
$\operatorname{Pic}_{X/S}^\circ$, it is unique. Its formation is consequently
compatible with base change.

\noindent(ii) For such an abelian subscheme to exist, it is enough that one
exist after every base change $S'\longrightarrow S$ with $S'$ local Artinian.
If $S$ is the spectrum of a local ring, it is enough to test the
$S'=\operatorname{Spec}(A_n)$, where $A_n=A/\mathfrak m^{n+1}$. If $S$ is
reduced, it is likewise enough to test the base changes
$S'\longrightarrow S$ in which $S'$ is the spectrum of a discrete valuation
ring (complete with algebraically closed residue field, if desired).

\noindent(iii) Suppose that $A$ exists, and let
$B=\operatorname{Alb}^0(X/S)$ be the dual abelian scheme (that is,
$B=\operatorname{Pic}_{A/S}^\circ$ [7]). Then one can canonically construct a
principal homogeneous space $P=\operatorname{Alb}^1(X/S)$ under $B$, together
with an $S$-morphism $X\longrightarrow P$ universal for $S$-morphisms from $X$
to para-abelian schemes (that is, principal homogeneous spaces under abelian
schemes). The formation of $\operatorname{Alb}^0(X/S)$,
$\operatorname{Alb}^1(X/S)$, and $X\longrightarrow\operatorname{Alb}^1(X/S)$
commutes with base change.

We sketch the proof.

Part (i) is a general rigidity property of abelian subschemes of commutative
group schemes: if two such subschemes coincide set-theoretically at a point
$s\in S$, they coincide over the entire connected component

\src{238}
of $s$ [7]. (This result generalizes a classical theorem of Chow.)

(ii) Hilbert schemes show that the functor associating to every $S'$ over $S$
the set (having either one or zero elements) of canonical abelian subschemes of
$(\operatorname{Pic}_{X/S})\times_S S'$ is represented by a scheme $T$ of
finite type over $S$. By (i), $T\longrightarrow S$ is a monomorphism, and by
(3.2) it is surjective. The existence of a canonical abelian subscheme of
$\operatorname{Pic}_{X/S}$ means that $T$ has a section over $S$, equivalently
that $T\longrightarrow S$ is an isomorphism. This is also equivalent to
$T\longrightarrow S$ being étale, or, when $S$ is reduced, to
$T\longrightarrow S$ being proper. This proves (ii).

(iii) Directly from the definition of $\operatorname{Pic}_{X/S}$, for every
abelian scheme $C$ over $S$, giving an $S$-morphism from $X$ to a principal
homogeneous space under $C$ is equivalent to giving a group homomorphism
$C'\longrightarrow\operatorname{Pic}_{X/S}$, where $C'$ is the dual abelian
scheme of $C$. If the canonical abelian subscheme $A$ of
$\operatorname{Pic}_{X/S}$ exists, these homomorphisms necessarily factor
through $A$ (as can also be seen by using points of finite order). This proves
(iii).

\result{REMARKS}{3.4} We shall denote by
$\operatorname{Pic}_{X/S}^{\circ\circ}$ the canonical abelian subscheme of
$\operatorname{Pic}_{X/S}$, when it exists. Unfortunately it does not always
exist, as is seen by varying the Igusa surface infinitesimally (through the
first-order modular deformation). It may nevertheless be that
$\operatorname{Pic}_{X/S}^{\circ\circ}$ exists at least when $S$ is reduced,
or equivalently, by (ii), when $S$ is the spectrum of a discrete valuation
ring. Let $X_0,X_1$ be the special and generic fibers of $X$, and put
\[
A_1=\operatorname{Pic}_{X_1/K}^\circ,
\]
where $K$ is the fraction field of the valuation ring $V$. By Koizumi's
theorem there is an essentially unique abelian scheme $A$ over $S$ whose
generic fiber is $A_1$. As in (2.1), (i), one sees readily (now assuming that
$X$ is smooth over $S$) that the identity morphism of $A_1$ extends to a
morphism
\[
A\longrightarrow\operatorname{Pic}_{X/S}.
\]
Hence there is a homomorphism
\[
\tag{*}A_0\longrightarrow\operatorname{Pic}_{X_0/k}^{\circ\circ},
\]
which is readily shown to be surjective with kernel a finite $p$-group, where
$p$ is the characteristic exponent of the residue field $k$

\src{239}
(again by using points of finite order). With this understood, the following
conditions on $X/S$ are equivalent:

\medskip
\noindent a. The preceding homomorphism (*) is an isomorphism (which Shimura
would express by saying that formation of the “Picard variety” is “compatible
with specialization”).

\noindent b. $\operatorname{Pic}_{X/S}^{\circ\circ}$ exists (in which case it
is none other than $A$).

\noindent c. (For reference) The
$\operatorname{Pic}_{X_n/S}^{\circ\circ}$ exist.

By the observation about the kernel of (*), condition (a) holds when the
residue characteristic is zero, but this result will be considerably
generalized in (3.5).

Of course, if $\operatorname{Pic}_{X/S}$ is smooth over $S$ at the points of
$\operatorname{Pic}_{X/S}^\circ$, then the latter is open in
$\operatorname{Pic}_{X/S}$ (cf. (1.7)). With the induced structure it is
therefore an abelian subscheme of $\operatorname{Pic}_{X/S}$, hence equal to
$\operatorname{Pic}_{X/S}^{\circ\circ}$, which exists in this case. But much
more is true:

\result{THEOREM}{3.5} Under the hypotheses of (3.3), let $s\in S$ be such that
$\operatorname{Pic}_{X_s/k(s)}$ is smooth over $k(s)$ (or, equivalently,
\[
\dim\operatorname{Pic}_{X_s/k(s)}=\dim H^1(X_s,\mathcal O_{X_s})).
\]
Then there is an open neighborhood $U$ of $s$ such that
$\operatorname{Pic}_{X/S}$ is smooth over $S$ at the points of
$\operatorname{Pic}_{X/S}^\circ|U$, which is consequently an open abelian
subscheme of $\operatorname{Pic}_{X/S}|U$. A fortiori,
$\operatorname{Pic}_{X|U/U}^{\circ\circ}$ exists.

We indicate the principle of the proof. The preceding discussion reduces us
to the case in which $S$ is the spectrum of a local Artinian ring $A$, and we
argue by induction on the infinitesimal order of $A$. We may therefore suppose
that $\operatorname{Pic}_{X_n/A_n}^\circ$ is smooth over $A_n$ and need only
prove that $\operatorname{Pic}_{X_{n+1}/A_{n+1}}^\circ$ is smooth over
$A_{n+1}$. For this it is enough to construct an abelian scheme $P_{n+1}$ over
$A_{n+1}$ extending
\[
P_n=\operatorname{Pic}_{X_n/A_n}^\circ,
\]
together with an invertible Module $\mathcal L_{n+1}$ on
$X_{n+1}\times_{A_{n+1}}P_{n+1}$ extending the invertible Module
$\mathcal L_n$ on $X_n\times_{A_n}P_n$ that occurs in the definition of the
Picard scheme $\operatorname{Pic}_{X_n/A_n}$ as the solution of a universal
problem. (N.B. --- We may suppose that $X$ has a section over $S$~...) The
construction uses the following key result: every abelian scheme defined over
a quotient of a local Artinian ring can be extended (in other words, the
absolute “formal moduli scheme” ([3], II) of an abelian scheme over an
algebraically closed field is smooth over the ring of

\src{240}
Witt vectors of $k$). This result itself follows from the general formal
properties of the obstruction to lifting and from the group operations.
Assuming it, first extend $P_n$ arbitrarily to $P_{n+1}$. One then encounters
an obstruction to lifting $\mathcal L_n$, lying in
\[
H^2(X_0\times P_0,\mathcal O_{X_0\times P_0})\otimes_k V
\qquad(\text{where }V=\mathfrak m^{n+1}/\mathfrak m^{n+2}),
\]
and in fact in the subspace
\[
H^1(X_0,\mathcal O_{X_0})\otimes
H^1(P_0,\mathcal O_{P_0})\otimes_k V
\]
(because the restriction of $\mathcal L_n$ to each of the two factors $X_n$
and $P_n$ is trivial). This last space is precisely
\[
H^1(P_0,\mathcal T_{P_0/k})\otimes V,
\]
where $\mathcal T_{P_0/k}$ is the tangent sheaf of $P_0/k$, and it is therefore
also the space measuring the indeterminacy in lifting $P_n$ to $P_{n+1}$ ([3],
II). Consequently this lift can be corrected (in one and only one way, as it
should be) so as to kill the obstruction to lifting $\mathcal L_n$.

\result{COROLLARY}{3.6} Under the hypotheses of (3.5),
$R^1f_*(\mathcal O_X)$ is a locally free Module on $S$ near $s$, and its
formation commutes with base change.

Indeed, this Module is the tangent Module of $\operatorname{Pic}_{X/S}$ along
the identity section.

\result{REMARK}{3.7} The same argument as in (3.5) shows that if $S'$ is a
subscheme of $S$ defined by a coherent nilpotent Ideal, and if one assumes only
that $\operatorname{Pic}_{X'/S'}^\circ$ exists and has smooth fibers, then
$\operatorname{Pic}_{X/S}^\circ$ necessarily exists and is an abelian scheme
over $S$. This permits Picard schemes to be constructed in some cases despite
the absence of a projectivity hypothesis; for example, the dual abelian scheme
of an abelian scheme over an Artinian ring always exists. On the other hand,
applying (3.6) when $X$ is an abelian scheme over $S$, and using the known
exterior-algebra structure of $H^*(X_s,\mathcal O_{X_s})$ over
$H^1(X_s,\mathcal O_{X_s})$ (Rosenlicht--Serre), one finds that for every $i$,
$R^if_*(\mathcal O_X)$ is locally free and, more precisely, isomorphic to the
$i$th exterior power of $R^1f_*(\mathcal O_X)$.

\result{REMARK}{3.8} For a projective smooth morphism
$f:X\longrightarrow S$, with $S$ reduced and all residue characteristics zero,
(3.6) was known by transcendental methods as a consequence of Hodge theory. In
fact, all the $R^pf_*(\Omega^q_{X/S})$ are then locally free. There are, on the
other hand, counterexamples

\src{241}
in mixed characteristic to the local freeness of $R^1f_*(\mathcal O_X)$,
obtained from Serre varieties ([8], p.~50). There seems to be no counterexample
in equal characteristic.

\result{COROLLARY}{3.8} Under the hypotheses of (3.5),
$\operatorname{Pic}_{X/S}|U$ is smooth over $U$ at every point of
$\operatorname{Pic}_{X/S}^\sigma|U$.

Apply (2.10), (i). In particular, taking (2.10), (iii), into account:

\result{COROLLARY}{3.9} Suppose that all residue characteristics of $S$ are
zero. Then $\operatorname{Pic}_{X/S}^\tau$ is smooth over $S$.

\result{REMARK}{3.10} If $\operatorname{Pic}_{X/S}^\tau$ is proper over $S$,
it follows, for example, that the torsion subgroup of the Néron--Severi groups
of the geometric fibers of $f$ is constant on each connected component of $S$
(as is also clear by transcendental methods when $f$ is smooth and projective).
A direct application of (2.5) shows more generally that if $\operatorname{Pic}_{X/S}^\tau$ is proper over $S$ (for example, if $f:X\longrightarrow S$ is smooth and projective) and $q$ is a prime distinct
from all residue characteristics of $S$, then the $q$-primary component of the
torsion subgroups of the Néron--Severi groups of the geometric fibers of $X/S$
is constant on each connected component of $S$. In characteristic $p>0$ this
is no longer true for the $p$-primary component of the torsion. It remains
possible, however, that the rank over $k(s)$ of
\[
\operatorname{Pic}_{X_s/k(s)}^\tau/
\operatorname{Pic}_{X_s/k(s)}^{\circ\circ}=T_{X_s/k(s)}
\]
is locally constant. When $S$ is reduced, one can show that this is equivalent
to the existence of $\operatorname{Pic}_{X/S}^{\circ\circ}$ together with
the flatness of $\operatorname{Pic}_{X/S}^\tau$ over $S$, and it is enough to
test the question when $S$ is the spectrum of a discrete valuation ring. This
is what I have verified in the few examples I have examined; but since the
corresponding assertion with $S$ Artinian is false (cf. Remarks (2.9) and
(3.4)), one should probably not indulge in excessive optimism.

\fsection{4}{The finiteness theorem for the Picard scheme}

Let $f:X\longrightarrow S$ be a projective flat morphism such that
$\operatorname{Pic}_{X/S}$ exists. Consideration of the “Hilbert polynomials”
$Q\in\mathbf Q[t]$ decomposes $\operatorname{Pic}_{X/S}$ into a disjoint union
of open subsets $\operatorname{Pic}_{X/S}^Q$. Without further hypotheses, for
example ones ensuring that $\operatorname{Pic}_{X/S}$ is separated over $S$,
these open subsets need not be of finite type over $S$; a counterexample is
obtained when $X$ is “a conic degenerating into two lines.” It is

\src{242}
possible, however, that they are of finite type when $f$ is separable with
geometrically irreducible fibers. The question is connected with whether
$\operatorname{Pic}_{X/S}^\tau$ is of finite type over $S$, which may be true
without any hypothesis on the fibers of $X/S$. When
$f:X\longrightarrow S$ is smooth, note that
$\operatorname{Pic}_{X/S}^\tau$ is contained in one of the
$\operatorname{Pic}_{X/S}^Q$ (because on a nonsingular projective variety
torsion equivalence is finer than—in fact, by Matsusaka, identical to—numerical
equivalence for divisors), and is therefore of finite type over $S$ if the
$\operatorname{Pic}_{X/S}^Q$ are. The finiteness questions just raised retain
an evident meaning even without assuming the existence of
$\operatorname{Pic}_{X/S}$: they say that certain families of invertible
Modules are “bounded” in the sense of [3], IV. For every algebraically closed
field $k$, consider the integral subschemes of $\mathbf P_k^r$ of dimension
$n$ and degree $d$, and the invertible Modules on such preschemes with Hilbert
polynomial $Q$, where $r,n,d,Q$ are fixed. The problem is to show that the
family of these Modules (viewed as coherent Modules on the $\mathbf P_k^r$) is
bounded, that is, can be parametrized by a scheme of finite type over
$\mathbf Z$.

Using Matsusaka's method [6], a rather technical proof (using the “equivalence
criteria” in a suitable form) gives an affirmative answer when one restricts
to nonsingular subvarieties of $\mathbf P_k^r$. More precisely:

\result{THEOREM}{4.1} Let $f:X\longrightarrow S$\footnote{The printed text
reads “Let: $X\longrightarrow S$”; the morphism used below is $f$.} be a smooth
projective morphism with geometrically connected fibers, with $S$ noetherian.
Let $\mathcal O_X(1)$ be a relatively very ample Module on $X$, and let $E$ be
a subset of $\operatorname{Pic}_{X/S}$ corresponding to a family
$(\mathcal L_i)$ of invertible Modules on the geometric fibers
$\overline X_{s_i}$ of $X/S$. Let $D_i$ be a divisor (not necessarily
effective) on $\overline X_{s_i}$ defining $\mathcal L_i$, let
\[
a_n^{(i)}t^n+\cdots+a_0^{(i)}
\]
be the Hilbert polynomial of $\mathcal L_i$, and let $\xi=\xi_i$ be a divisor
defining $\mathcal O_X(1)$, that is, a hyperplane section. The following
conditions are equivalent:

\noindent a. $E$\footnote{The printed text reads “$Q$.” The subset introduced
in the statement, whose quasi-compactness expresses boundedness of the family,
is $E$.} is quasi-compact, that is, contained in an open subset of finite type
over $S$; equivalently, the family $(\mathcal L_i)$ is bounded.

\noindent b. $E$ is contained in the union of finitely many of the sets
$\operatorname{Pic}_{X/S}^Q$, $Q\in\mathbf Q[t]$; equivalently, only finitely
many Hilbert polynomials occur among the $\mathcal L_i$.

\noindent b'. (If all fibers of $X/S$ have the same dimension $n$.) The two
coefficients $a_{n-1}^{(i)}$ and $a_{n-2}^{(i)}$ of the Hilbert polynomials of
the $\mathcal L_i$ range over finite sets.

\src{243}

\result{COROLLARY}{4.2} Let $f:X\longrightarrow S$ be a smooth projective
morphism with geometrically connected fibers. Then the schemes
$\operatorname{Pic}_{X/S}^Q$ and $\operatorname{Pic}_{X/S}^\tau$
($Q\in\mathbf Q[t]$) are projective over $S$.

They are of finite type over $S$ by (4.1), so (2.4) applies.

\section*{BIBLIOGRAPHY}
\addcontentsline{toc}{section}{Bibliography}

\begin{description}
\item[{[1]}] CHEVALLEY (Claude). -- Sur la théorie de Picard,
\emph{Amer. J. of Math.}, vol. 82, 1960, pp. 435--490.

\item[{[2]}] GROTHENDIECK (A.) and DIEUDONNÉ (J.). -- Éléments de géométrie
algébrique, Chapter I and following. -- Paris, Presses universitaires de
France, 1960--1961 (Institut des Hautes Études Scientifiques, 4, 8, 11,
\ldots).

\item[{[3]}] GROTHENDIECK (Alexander). -- Technique de descente et théorèmes
d'existence en géométrie algébrique, I--V, \emph{Séminaire Bourbaki}, vol. 12,
1959/60, no. 190, 29 pp. and no. 195, 22 pp.; vol. 13, 1960/61, no. 212,
20 pp. and no. 221, 28 pp.; vol. 14, 1961/62, no. 232, 19 pp.

\item[{[4]}] GROTHENDIECK (Alexander). -- Séminaire de géométrie algébrique,
1960/61. -- Paris, Institut des Hautes Études Scientifiques, 1961.

\item[{[5]}] GROTHENDIECK (Alexander). -- Séminaire de géométrie algébrique,
1961/62, edited by a group of auditors. -- Paris, Institut des Hautes Études
Scientifiques (forthcoming).

\item[{[6]}] MATSUSAKA (T.). -- The criteria for algebraic equivalence and the
torsion group, \emph{Amer. J. Math.}, vol. 79, 1957, pp. 53--66.

\item[{[7]}] MUMFORD (D.) and TATE (J.). -- Séminaire de géométrie algébrique,
Harvard University, Spring Term 1962 (forthcoming).

\item[{[8]}] SERRE (Jean-Pierre). -- Sur la topologie des variétés algébriques
en caractéristique $p$, \emph{Symposium internacional de topología algebraica}
[1956, Mexico], pp. 24--53. -- Mexico, Universidad nacional autónoma, 1958.
\end{description}

\par\smallskip
\begingroup\footnotesize
\noindent\textit{through 0058. In 3.2 the print assumes $X$ normal and proper over $k$, without geometric normality (C122). In 3.10 both printed properness premises concern the whole Picard scheme, not $\operatorname{Pic}_{X/S}^{\tau}$ (C124). C123 restores the printed $\sigma$ in 2.10(i) and 3.8: the circle and $\tau$ were transcription errors, not printed readings.}\par
\endgroup
